\documentclass[11pt]{article}
\usepackage[margin=1in]{geometry}
\usepackage{amsmath,amssymb,amsthm,mathtools}
\usepackage{graphicx}
\usepackage{microtype}
\usepackage[hidelinks]{hyperref}

\newtheorem{theorem}{Theorem}
\newtheorem{lemma}{Lemma}
\newtheorem{remark}{Remark}
\newcommand{\C}{\mathbb C}
\newcommand{\Hh}{\mathcal H}
\newcommand{\LO}{\mathrm{LO}}
\newcommand{\LOCC}{\mathrm{LOCC}}
\newcommand{\SEP}{\mathrm{SEP}}
\newcommand{\tr}{\operatorname{Tr}}

\title{An Unbounded Oriented One-Way LOCC Gap by Subsystem Swap}
\author{Candidate full answer to the first question in Section 7.3 of arXiv:1406.1959}
\date{26 August 2026}

\begin{document}
\maketitle

\begin{center}
\fbox{\parbox{0.91\textwidth}{\textbf{Status.}
Candidate full result, likely valid, pending human review.  For the oriented
one-way class $\LOCC^{A\to B}$ displayed in the source, there are bipartite
state pairs whose $\LOCC$ distinguishability norm equals $2$ while their
$\LOCC^{A\to B}$ norm is $O(d^{-1/2})$.  Thus no dimension-free comparison
$\|\cdot\|_{\LOCC^{A\to B}}\geq c\|\cdot\|_{\LOCC}$ is possible.}}
\end{center}

\section{Source question and convention}

Aubrun and Lancien define the oriented one-way class
$\LOCC^{A\to B}$ by POVMs of the form
\[
 \bigl(M_i\otimes N_{i,j}\bigr)_{i,j},\qquad
 \sum_iM_i=I_A,\quad \sum_jN_{i,j}=I_B,
\]
so Alice measures first and Bob's measurement may depend on Alice's outcome
\cite[Section 2.3]{AL}.  Section 7.3 asks whether the gap between this class
and unrestricted finite-round $\LOCC$ can be unbounded, or whether an
absolute $c>0$ satisfies
\[
                  \|\cdot\|_{\LOCC^{A\to B}}
                  \geq c\|\cdot\|_{\LOCC}                 \tag{1}
\]
in every local dimension.

\begin{figure}[ht]
\centering
\includegraphics[page=20,width=0.93\textwidth,trim=55 125 55 420,clip]{source_paper.pdf}
\caption{The printed question in Section 7.3 of the source, arXiv PDF page 20.}
\end{figure}

The direction matters.  We write $\LOCC^{B\to A}$ for the analogous class
with Bob measuring first.  The result below answers the question for the
oriented class exactly as displayed.  It does not address a different
convention in which ``one-way LOCC'' means the union of the two orientations.

For a POVM $\mathsf M=(M_a)_a$ and a Hermitian operator $\Delta$, recall
\[
                  \|\Delta\|_{\mathsf M}
                  =\sum_a|\tr(\Delta M_a)|,
\]
and the norm of a measurement class is the supremum over its POVMs.

\section{Reverse communication collapses on classical-quantum inputs}

\begin{lemma}[Classical-side collapse]\label{lem:collapse}
Let
\[
        \Delta=\sum_{i=1}^d |i\rangle\!\langle i|\otimes\Delta_i
        \in\Hh(\C^d\otimes\C^d),                         \tag{2}
\]
where every $\Delta_i$ is Hermitian.  Then
\[
 \|\Delta\|_{\LOCC^{B\to A}}
 =\|\Delta\|_{\LO}
 =\sup_{\mathsf N=(N_j)_j}\sum_{i=1}^d
                         \|\Delta_i\|_{\mathsf N}.        \tag{3}
\]
\end{lemma}

\begin{proof}
The last expression is attained within $\LO$: Alice measures in the basis
$(|i\rangle\!\langle i|)_i$ and Bob uses $\mathsf N$.  Conversely, consider
an arbitrary $B\to A$ one-way POVM.  Its effects have the form
\[
                         M_{j,k}\otimes N_j,
\]
where $\mathsf N=(N_j)_j$ is Bob's first measurement and, for every $j$,
$(M_{j,k})_k$ is Alice's conditional POVM.  Put
$m_{j,k,i}=\langle i|M_{j,k}|i\rangle$.  These numbers are nonnegative and
$\sum_km_{j,k,i}=1$.  Hence
\begin{align*}
 \sum_{j,k}\left|\tr\bigl(\Delta(M_{j,k}\otimes N_j)\bigr)\right|
 &=\sum_{j,k}\left|\sum_i m_{j,k,i}\tr(\Delta_iN_j)\right|\\
 &\leq\sum_{j,k,i}m_{j,k,i}|\tr(\Delta_iN_j)|\\
 &=\sum_{i,j}|\tr(\Delta_iN_j)|
   =\sum_i\|\Delta_i\|_{\mathsf N}.
\end{align*}
Taking the supremum gives
$\|\Delta\|_{\LOCC^{B\to A}}\leq\|\Delta\|_{\LO}$, while the reverse
inequality follows from $\LO\subset\LOCC^{B\to A}$.
\end{proof}

The lemma says that communication arriving at a party whose register is
already classical cannot improve the norm: any conditional processing of
that classical register is dominated by reading it outright at the start.

\section{The swapped data-locking pair}

The source proves the following theorem.  There is a universal constant $C$
such that, for every local dimension $d$, one can find states $\rho_d,\sigma_d$
on $\C^d\otimes\C^d$ with
\[
 \|\rho_d-\sigma_d\|_{\LOCC^{A\to B}}=2,
 \qquad
 \|\rho_d-\sigma_d\|_{\LO}\leq\frac{C}{\sqrt d}.        \tag{4}
\]
Moreover, their construction is classical on Alice's side:
\[
 \rho_d=\frac1d\sum_{i=1}^d|i\rangle\!\langle i|\otimes\rho_i,
 \qquad
 \sigma_d=\frac1d\sum_{i=1}^d|i\rangle\!\langle i|\otimes\sigma_i,       \tag{5}
\]
where $\rho_i$ and $\sigma_i$ have orthogonal supports \cite[Theorem 3]{AL}.

\begin{figure}[ht]
\centering
\includegraphics[page=7,width=0.93\textwidth,trim=55 315 55 185,clip]{source_paper.pdf}
\caption{Theorem 3 and the classical-quantum construction in the source,
arXiv PDF page 7.}
\end{figure}

Let $S(x\otimes y)=y\otimes x$ be the tensor-factor swap and define
\[
             \widetilde\rho_d=S\rho_dS^*,\qquad
             \widetilde\sigma_d=S\sigma_dS^*.
\]

\begin{theorem}[Unbounded oriented one-way gap]\label{thm:main}
There is a universal constant $C$ such that
\[
 \|\widetilde\rho_d-\widetilde\sigma_d\|_{\LOCC}=2,
 \qquad
 \|\widetilde\rho_d-\widetilde\sigma_d\|_{\LOCC^{A\to B}}
     \leq\frac{C}{\sqrt d}.                               \tag{6}
\]
Consequently
\[
 \sup_{0\neq\Delta\in\Hh(\C^d\otimes\C^d)}
 \frac{\|\Delta\|_{\LOCC}}
      {\|\Delta\|_{\LOCC^{A\to B}}}
 \geq \frac{2}{C}\sqrt d,
\]
and the dimension-free inequality \emph{(1)} is false.
\end{theorem}

\begin{proof}
Swapping the factors converts an $A\to B$ protocol into a $B\to A$ protocol
and conversely.  Therefore Lemma \ref{lem:collapse}, applied to
$\rho_d-\sigma_d$ in (5), and (4) give
\[
 \|\widetilde\rho_d-\widetilde\sigma_d\|_{\LOCC^{A\to B}}
 =\|\rho_d-\sigma_d\|_{\LOCC^{B\to A}}
 =\|\rho_d-\sigma_d\|_{\LO}
 \leq\frac{C}{\sqrt d}.
\]
Unrestricted $\LOCC$ is invariant under exchanging the parties.  The
$A\to B$ protocol in (4) distinguishes $\rho_d$ and $\sigma_d$ perfectly, so
after the swap the corresponding $B\to A$ protocol is an unrestricted
$\LOCC$ protocol that distinguishes $\widetilde\rho_d$ and
$\widetilde\sigma_d$ perfectly.  Its norm is therefore $2$, the maximum
possible for a difference of states.
\end{proof}

\section{Scope, verification, and literature boundary}

The proof uses no new probabilistic estimate: all asymptotics come from the
source's Theorem 3.  The new step is the exact collapse lemma, followed by
the tensor-factor swap.  The companion script
\texttt{code/verify\_collapse.py} tests the triangle inequality behind the
lemma on randomly generated classical-quantum operators and reverse one-way
POVMs; the proof itself is the displayed scalar calculation.

Bounded exact-question, exact-title, keyword, and citation searches through
26 August 2026 found work separating one-way and two-way protocols in other
hypothesis-testing regimes, but no later primary source recording this
$\Omega(\sqrt d)$ norm separation or observing the subsystem-swap consequence
of Theorem 3.  The result should be treated as a candidate full answer under
the source's displayed oriented convention, pending expert review.

\begin{remark}[Convention boundary]
If $\LOCC^{1\text{-way}}$ is instead defined as
$\LOCC^{A\to B}\cup\LOCC^{B\to A}$, the swapped pair is perfectly
distinguished by its $B\to A$ member and gives no separation.  The stronger
union-of-directions question remains open here.
\end{remark}

\section{\hspace{0.25em}Human review recommendation}

Check three points: the source's arrow is oriented exactly as in its displayed
definition; the absolute-value triangle inequality in Lemma
\ref{lem:collapse} permits no adaptive advantage on the classical register;
and tensor swap preserves unrestricted $\LOCC$ while reversing the one-way
orientation.  These checks make (6) immediate from the source's Theorem 3.

\begin{thebibliography}{9}
\bibitem{AL}
G. Aubrun and C. Lancien,
\emph{Locally restricted measurements on a multipartite quantum system:
data hiding is generic}, Quantum Information and Computation
\textbf{15} (2015), 513--540; arXiv:1406.1959.
\end{thebibliography}

\end{document}
